\documentclass[11pt]{article}
\usepackage[margin=1in]{geometry}
\usepackage{amsmath,amssymb,amsfonts,mathtools,bm}
\usepackage[T1]{fontenc}
\usepackage[utf8]{inputenc}
\title{Inline and Display Mathematics}
\author{Source-backed grouped formula sample}
\date{}
\begin{document}
\maketitle
\section{Expressions}
\paragraph{Expression 1.} The following expression is taken from a source-backed formula corpus.

\begin{align*}S = -\frac{1}{2}\int\!d\sigma d\tau[\partial_\mu\chi_1 \partial^\mu\chi_1 + 2\chi_2\partial_\mu\chi_1 \partial^\mu\chi_3 + \frac{1}{1+{\chi_1}^2}\partial_\mu\chi_2 \partial^\mu\chi_2 + ({\chi_2}^2+\frac{1}{1+{\chi_1}^2})\partial_\mu\chi_3 \partial^\mu\chi_3 +\partial_\mu t \partial^\mu t].\end{align*}

\paragraph{Expression 2.} The following expression is taken from a source-backed formula corpus.

\begin{align*}\left({}^{c}D_{0+}^{\alpha,\alpha',\beta,\beta',\gamma}t^{\rho-1}\right)(x)=\frac{\Gamma(\rho)\Gamma(\alpha-\beta+\rho-m)\Gamma(\alpha+\alpha'+\beta'-\gamma+\rho-m)}{\Gamma(-\beta+\rho-m)\Gamma(\alpha+\alpha'-\gamma+\rho)\Gamma(\alpha+\beta'-\gamma+\rho-m)}x^{\alpha+\alpha'-\gamma+\rho-1}.\end{align*}

\paragraph{Expression 3.} The following expression is taken from a source-backed formula corpus.

\begin{align*} \phi ( {\vec {X}, \varepsilon} ) = \dot { \vec{X} } \cdot ( {\ddot {\vec {X}} \wedge \dddot{\vec {X}}\wedge \ldots\wedge \mathop {\vec {X}}\limits^{\left( n \right)} } ) = \det( {\dot {\vec {X}},\ddot {\vec {X}}, \dddot{\vec {X}},\ldots ,\mathop{\vec {X}}\limits^{\left( n \right)} } ) = 0\end{align*}

\paragraph{Expression 4.} The following expression is taken from a source-backed formula corpus.

\begin{align*}\vert p,\sigma,q,\lambda \colon f(p) \rangle = exp\biggl[-e\int\sum_{\lambda=1,2} {{dk^+} \over {\sqrt{2k^+}}} \int {{d^2k_\perp} \over{(2\pi)^{3 \over 2}}}\bigl[f(k,\lambda \colon p)a^\dagger(k,\lambda) - f^*(k,\lambda \colon p)a(k,\lambda)\bigr] \biggr]\vert p,\sigma,q,\lambda \rangle\;.\end{align*}

\paragraph{Expression 5.} The following expression is taken from a source-backed formula corpus.

\begin{align*}\Psi_2 \left( \rho_1(y_1, y_1^\prime)^{p} , \boldsymbol{d}_{\mathcal{S}_2}(s_2, s_2^\prime)^p \right) = \Psi_2 \circ \psi_2 \left( \boldsymbol{d}_{\mathcal{Y}_1}(y_1, y_1^\prime)^{p} , \boldsymbol{d}_{\mathcal{Y}_2}(y_2, y_2^\prime )^p , \boldsymbol{d}_{\mathcal{X}}(x, x^\prime)^p \right).\end{align*}
\end{document}
